%% Unit 2 -- Driving opencode: the commands and keys worth knowing.
%% Commands are grouped by theme, several to a slide, as a two-column booktabs
%% table: how to run it on the left, what it does on the right. Default keys are
%% opencode's out-of-the-box bindings; the leader is ctrl+x.
%% No preamble here; \input by the master (disciplined_vibe.tex).

% The command-table helpers (\gcmd, \gkey, \gpal, \gsep and the cmdtable
% environment) now live in the master preamble (disciplined_vibe.tex), since the
% command tables are shared across units 2, 3, and 4.

\sectionsubtitle{When Language becomes Action}
% ======================================================================
\section[Vibing]{Driving opencode}
% ======================================================================



\iffalse
% -------------------------------------------------------- historical arc
\begin{frame}{A short history of NLP}
  \begin{columns}[T,onlytextwidth]
    \begin{column}{0.24\textwidth}
      \textbf{Pre-2017}\par\smallskip
      {\small N-gram and statistical models, and word embeddings
      such as Word2Vec and GloVe. Early
      chatbots existed, such as ELIZA.}
    \end{column}
    \begin{column}{0.24\textwidth}
      \textbf{2017: the Transformer}\par\smallskip
      {\small The \alert{Transformer} enables highly
      efficient, scalable natural language processing (NLP).}
    \end{column}
    \begin{column}{0.24\textwidth}
      \textbf{2018 to 2021: LLMs}\par\smallskip
      {\small Billions of parameters. Fast progress on pre-training, few-shot
      learning, and RLHF yields BERT, GPT-2, T5, LLaMA.}
    \end{column}
    \begin{column}{0.24\textwidth}
      \textbf{2022, or 0 A.ChatGPT}\par\smallskip
      {\small Mainstream conversational LLMs: GPT-3, and
      from there into vision, language, action models such as GR00T
      and SmolVLA.}
    \end{column}
  \end{columns}
\end{frame}
\fi

% -------------------------------------------------------- tool injection
\begin{frame}{Designing Interfaces by Injection}
  \vskip0.4em
  \begin{withmargin}
    \begin{tdmain}
      The model only reads and writes tokens. A tool still needs an
      interface, so we make one the same way as before: by injecting its
      description into the prompt, then reading the call back out of the reply.
      \vskip0.8em
      \begin{itemize}
        \item The \alertbf{system prompt} specifies the tool, its arguments, and the
          format for calling it.
        \item The model writes that call as ordinary text; a \alertbf{parser} spots
          it and runs the real tool.
        \item The result comes back as tokens and get injected as observations, and the model continues.
      \end{itemize}
    \end{tdmain}
    \begin{tdmargin}
      \begin{tdcodet}{System prompt}
        tool weather(city)\\
        call <tool>...</tool>
      \end{tdcodet}
      \vskip6pt
      \begin{tdcodet}{LLM Tool use}
        <tool>\\
        weather(Ravensburg)\\
        </tool>
      \end{tdcodet}
    \end{tdmargin}
  \end{withmargin}
\end{frame}


% ----------------------------------------------------- what MCP is
\begin{frame}{MCP and plug in outside tools}
  \begin{withmargin}
    \begin{tdmain}
      The \alertbf{Model Context Protocol} is a standard way to connect the agent
      to outside tools and data: a database, a web service, a file store, a program, each
      exposed through a small server.
      \vskip1.0em
      \begin{tdblock}
        MCP is a
        protocol. The agent host connects, discovers them at runtime, and calls
        them. Write to the standard once, and any MCP client can use it.       
      \end{tdblock}
      \vskip0.5em
      {\small }
    \end{tdmain}
    \begin{tdmargin}
      This is the structured answer to agent-ready APIs. The engineering way with proper auth and 
      hardened.
    \end{tdmargin}
  \end{withmargin}
\end{frame}



% ----------------------------------------------------------------- action
\begin{frame}{The one Acting}
  \slidesources{nvidia2025gr00t,shukor2025smolvla}
  \sideimagecover[0.33]{amznghand}
  \begin{sidebody}
    \begin{tdblock}\small
      {\alert{\textbf{Agentic}}, from Latin \textit{agere}, ``to act,'' the
      root of \textit{agent}, one who acts.}
    \end{tdblock}
    \vskip0.15em
    {\footnotesize
    \begin{itemize}\itemsep1pt
      \item Emits its action in a fixed structured format using tags.
      \item Stops at \texttt{</action> or </tool>} once emitted.
      \item A parser then runs the action.
    \end{itemize}}
    \vskip0.15em
    {\footnotesize An action can be many things. Communication and subagent orchestration.
    Information gathering through search, queries,
    and retrieval. It uses tools, from API calls to code execution. Or it acts on
    the environment, driving digital or embodied interfaces.}
  \end{sidebody}
\end{frame}

% ============== The harness that wraps the agent ======================
% With thought, action, and observation on the table, opencode is the harness
% around that agent; sandboxing is the boundary it runs in.

% ----------------------------------------------------- what a harness is
\begin{frame}{The Harness}
  \slidesources{kumar2025harnesses}
  \sideimagecover[0.33]{harness}
  \begin{sidebody}
    The machine got an interface, now it's time for the human to get one to take
    control.
    \vskip0.6em
    \begin{tdblock}
      The \alertbf{harness} is everything wrapped around the raw model to make it
      reliable: the client, the system prompt, the skills, the tests, the
      permission gates, the loop that runs them.
    \end{tdblock}
    \vskip0.6em
    {\texttt{opencode} is your primary harness interface: type
    \alertbf{\texttt{/}} to open the command menu and pick by name.}
    \vskip0.5em
    {\small It has a set of \alertbf{built-in tools} which can be
    called directly: \texttt{bash}, \texttt{read}, \texttt{write},
    \texttt{edit}, \texttt{glob}, \texttt{grep}, \texttt{webfetch}.}
  \end{sidebody}
\end{frame}


% ====================== Getting oriented ==============================
\begin{frame}{Getting started}
  \begin{cmdtable}
    \gcmd{/init} &
      Scans the project and writes an \texttt{AGENTS.md} of your stack, scripts,
      and conventions, so every later session starts oriented and in \alertbf{context}. Run once per repo.\\[4pt]
    \gcmd{/help} &
      Opens the in-app help: the command list, the current keybinds, and where
      the docs live. Your first stop when a key surprises you.\\[4pt]
    \gcmd{/exit}\gsep\gkey{ctrl+x q} &
      Ends the session and drops you back to the shell. \gkey{ctrl+c} on an empty
      prompt quits too.\\
  \end{cmdtable}
\end{frame}



% ====================== Sessions ======================================
\begin{frame}{Steer by Context}
  \begin{cmdtable}
    \gcmd{/compact}\gsep\gkey{ctrl+x c} &
    Summarises the conversation into a short brief and continues from it,
    reclaiming context-window space without losing the thread. Don't wait too long.\\[4pt]
    \gcmd{/new}\gsep\gkey{ctrl+x n} &
      Starts a fresh session with an empty context. Your project stays; the
      conversation so far is dropped. Reach for it when the topic changes.\\[2pt]
    \gcmd{/sessions}\gsep\gkey{ctrl+x l} &
      Lists your earlier sessions and jumps to the one you pick.\\[2pt]
    \gkey{ctrl+r} &
      \textbf{Rename} the current session to a memorable title, so you find it
      again in the list.\\[2pt]
    \gkey{esc} &
      \textbf{Interrupt} the agent mid-run the moment it heads the wrong way,
      instead of waiting for it to finish.\\
  \end{cmdtable}
\end{frame}


% -------------------------------------------- gate vs cage (two layers)
\begin{frame}{Gating and sandboxing}
  \vskip0.3em
  {\small Gating keeps the human on the loop, the sanbox stops agent escapes. 
  Latest when you auto, you want a sandbox.
  }
  \vskip0.8em
  \begin{columns}[T,onlytextwidth]
    \begin{column}{0.47\textwidth}
      \begin{tdblock}
        \small
        \textbf{Gating}\par\smallskip
        \alertbf{opencode asks} before it edits or runs: \texttt{allow},
        \texttt{ask}, \texttt{deny}, set per tool in \texttt{opencode.json}.
        Ships in the box.
        \vskip0.5em
        Stops surprises. It cannot stop a command you already approved.
      \end{tdblock}
    \end{column}
    \begin{column}{0.47\textwidth}
      \begin{tdblock}
        \small
        \textbf{Sandboxing}\par\smallskip
        A container or VM. A project mounted, network restricted, \texttt{opencode}
        running inside it. You build it yourself.
        \vskip0.5em
        \alertbf{Stops escapes}. If something goes wrong it wrecks the box,
        not your machine.
      \end{tdblock}
    \end{column}
  \end{columns}
  \vskip0.9em
  \centering
\end{frame}
















% -------------------------------------------------------------- reasoning
% Two ways to get an inner monologue: elicit it by prompting (CoT, and its
% agentic extension ReAct) on the left, or train it in with RL (R1, o1) right.
\begin{frame}{Reasoning or the Inner Monologue}
  \slidesources{wei2022cot,yao2022react,deepseek2024r1}
  \vskip0.4em
  \begin{columns}[T,onlytextwidth]
    \begin{column}{0.48\textwidth}\small
      \textbf{Prompted at inference time}\par\smallskip
      \textit{reasoning coaxed out by the prompt, no weights changed.}
      \vskip0.5em
      \alertbf{Chain-of-thought} prompting appends a cue like \textit{``Let's think
      step by step,''} steering decoding toward a plan rather than a snap answer,
      so the model breaks the task into sub-tasks.
      \vskip0.5em
      We can interleave those reasoning steps with
      actions and observations so thinking and tool calls advance together, called ReAct.
    \end{column}
    \begin{column}{0.48\textwidth}\small
      \textbf{Trained into the weights}\par\smallskip
      \textit{reasoning trained not nudged.}
      \vskip0.5em
      Models like DeepSeek-R1 and OpenAI's o1 are trained with \alertbf{human feedback reinforcement
      learning} to reason before they answer.
      \vskip0.5em
      \begin{tdblock}
        R1 emits a thinking section between the tokens \texttt{<think>} and
        \texttt{</think>} before its answer.
      \end{tdblock}
    \end{column}
  \end{columns}
\end{frame}


% ====================== Interface and input ===========================
\begin{frame}{Interface and input}
  \begin{cmdtable}
    \gcmd{/editor}\gsep\gkey{ctrl+x e} &
      Opens your \texttt{\$EDITOR} to compose a long message with full editing,
      then sends it back to the prompt.\\[4pt]
    \gkey{ctrl+v} &
      \textbf{Paste} clipboard content into the prompt.\\[4pt]
    \gkey{ctrl+c} &
      \textbf{Clear} the prompt box in one stroke.\\[4pt]
    \gcmd{/details} &
      Toggles a tool call's full input and output, to inspect what
      ran in detail.\\[4pt]
    \gcmd{/thinking} &
      Toggles the model's reasoning \texttt{<think>} blocks, if the model has this capability.\\
  \end{cmdtable}
\end{frame}

% ====================== The surfaces you are driving ==================
% Four awareness frames: what a skill, MCP, harness, and sandbox actually
% are. We use only skills in this course; the other three are named so the
% students know them, and so Unit 3 can treat skills as verbs it invokes.

% (Skills frame moved to Unit 3, placed before the grill-me chain.)





% ------------------------------------------------------------- quest: vibe it
\begin{frame}{\alertbf{Vibe} an infinitely running banana}
  \sideimagecover[0.33]{banana_1}
  \sidecaption{Screenshot of a banana again.}
  \begin{sidebody}
    Now vibe an endless runner in a single \texttt{index.html}. 
    Then ask the LLM to run it in your browser 
    and play, then iterate.
    \begin{itemize}
      \item How easy is it to get one specific thing done?
      \item Watch your context and pay attention to the impact of a saturated context.
      \item Reflect on how productive this feels versus how productive you are.
    \end{itemize}
    \vspace{0.5em}
    \begin{tdblock}
      Feel the \alert{floor rising} already? Lines of code have never been cheaper.
    \end{tdblock}
  \end{sidebody}
\end{frame}

% References are collected in one shared package at the end of the deck
% (see disciplined_vibe.tex), not per unit.